% 2
\begin{frame}{오늘의 학습 목표}
  강의가 끝나면 다음을 할 수 있다.
  \begin{enumerate}
    \item \keyterm{알고리즘}, \keyterm{프로그램}, \keyterm{의사코드}를 구분한다.
    \item Maximum, Linear Search, Binary Search를 추적하고 설명한다.
    \item 입력 크기 $n$에 대한 기본 연산 수를 센다.
    \item $\Oh$, $\Om$, $\Th$의 의미를 수식과 말로 표현한다.
  \end{enumerate}
  \vfill
  \begin{block}{강의를 관통하는 질문}
    답만 맞는 절차를 넘어, \alert{왜 맞고 얼마나 잘 확장되는가?}
  \end{block}
\end{frame}

\section{Part A. 왜 알고리즘인가? (Why Algorithms Matter)}

% 3
\begin{frame}{같은 문제, 다른 선택}
  \begin{columns}[T,onlytextwidth]
    \column{.53\textwidth}
      \textbf{67센트를 거슬러 주자}\par
      동전: $25,10,5,1$센트
      \begin{center}
      \begin{tikzpicture}[node distance=2mm]
        \node[active] (a) {25};
        \node[active,right=of a] (b) {25};
        \node[active,right=of b] (c) {10};
        \node[active,right=of c] (d) {5};
        \node[active,right=of d] (e) {1};
        \node[active,right=of e] (f) {1};
      \end{tikzpicture}
      \end{center}
      큰 동전부터 선택하는 greedy 전략은 여기서는 6개를 쓴다.
    \column{.43\textwidth}
      \begin{alertblock}{하지만 질문이 남는다}
        \begin{itemize}
          \item 항상 답을 만드는가?
          \item 항상 최소 개수인가?
          \item 동전 체계가 바뀌어도 그런가?
        \end{itemize}
      \end{alertblock}
      최적임은 \textbf{증명}, 최적이 아님은 \textbf{반례}로 판단한다.
  \end{columns}
\end{frame}

\begin{frame}{탐욕법(Greedy)은 항상 최적인가?}
  \begin{columns}[c,onlytextwidth]
    \column{.48\textwidth}
      \begin{block}{동전 $\{1,3,4\}$, 목표 6}
        큰 동전부터 선택하면
        \[
          4+1+1=6 \qquad \textbf{3 coins}
        \]
      \end{block}
    \column{.48\textwidth}
      \begin{alertblock}{더 좋은 해}
        \[
          3+3=6 \qquad \textbf{2 coins}
        \]
      \end{alertblock}
  \end{columns}
  \vfill
  \begin{takeaway}
    그럴듯한 알고리즘이 항상 정확하거나 최적이지는 않다. \textbf{최적성은 증명하거나, 반례로 깨야 한다.}
  \end{takeaway}
\end{frame}

% 4
\begin{frame}{규모가 답을 바꾼다}
  \begin{columns}[c,onlytextwidth]
    \column{.52\textwidth}
      \begin{tikzpicture}[scale=1.05,transform shape]
        \draw[->,AlgoGray] (0,0)--(5.3,0) node[right] {$n$};
        \draw[->,AlgoGray] (0,0)--(0,3.2) node[above] {work};
        \draw[AlgoBlue,very thick] plot[smooth] coordinates {(0,0.2)(1,.55)(2,.9)(3,1.25)(4,1.6)(5,1.95)};
        \draw[AlgoOrange,very thick] plot[smooth] coordinates {(0,.15)(1,.3)(2,.65)(3,1.25)(4,2.15)(5,3.25)};
        \node[AlgoBlue] at (4.4,1.35) {linear};
        \node[AlgoOrangeText] at (4.4,2.65) {quadratic};
        \node[note,anchor=west] at (.15,3.0) {개념도: 축의 수치는 생략};
      \end{tikzpicture}
    \column{.44\textwidth}
      작은 입력에서는 구현 상수의 차이가 커 보인다. 하지만 $n$이 커지면 \alert{성장률}이 지배한다.
      \begin{takeaway}
        알고리즘은 실행 가능성뿐 아니라 \textbf{정확성(correctness)}과
        \textbf{확장성(scalability)}을 함께 결정한다.
      \end{takeaway}
  \end{columns}
  \vfill
  \muted{알고리즘이 왜 중요한지는 보았다. 이제 프로그램과 구분하고, 무엇을 알고리즘이라고 부를지 정리하자.}
\end{frame}
